\subsection*{IV 20: the opening contents and Commentaires (cote 25, sites 2--3)}
The opening contents list sixteen proposed sections, from \emph{Notations} to \emph{Traductions en langage de systèmes linéaires}. In the \emph{Commentaires}, Grothendieck describes the present \emph{Etat 0} as an assembly of detailed drafts for a later redaction, and expressly calls for a new plan. He names the current nos.\ 11, 12, 14 and 15 as material to be moved earlier. The continuation on site 3 places proposition 10.3 among results needing an earlier home and mentions further material on \emph{applications birationelles}. He proposes splitting the current paragraph 20, with the developments on grassmannians making one of the resulting parts. The projected no.\ 16 remains unwritten at the moment of these comments. These are changes of organization stated on the pages themselves; the edition preserves the successive headings and references. The doubt retained in the upper-right pencil mark of site 3 concerns a folio annotation; it does not supply a new section number. Sources: \href{https://grothendieck.umontpellier.fr/web/viewer.html?file=../25.pdf\#page=3}{cote 25, site 2} and \href{https://grothendieck.umontpellier.fr/web/viewer.html?file=../25.pdf\#page=4}{site 3}.
